\documentclass[12pt, a4paper, oneside]{ctexart}
\usepackage{amsmath, amsthm, amssymb, bm, xcolor, tikz, framed, graphicx, hyperref, mathrsfs}
\usepackage{bookmark}
\usepackage{pgfplots}
\pgfplotsset{compat=1.18}

\newcommand{\bra}[1]{\langle #1 |}
\newcommand{\ket}[1]{| #1 \rangle}
\newcommand{\braket}[2]{\langle #1 | #2 \rangle}

\title{Quantum Mechanics Homework (Week 14)}
\author{Tang Yi \\ [学号已脱敏]}
\usepackage[a4paper, left=2cm, right=2cm, top=2cm, bottom=2cm]{geometry}
\linespread{1.3} % 行距更紧凑
\definecolor{shadecolor}{RGB}{241, 241, 255}

\newcounter{problemname}
\newenvironment{problem}
  {\begin{shaded}\stepcounter{problemname}\par\noindent\textbf{题目 \theproblemname. }}
  {\end{shaded}\par}
\newenvironment{solution}
  {\par\noindent\textbf{题目 \theproblemname 的解答. }}
  {\par}
\newenvironment{note}
  {\par\noindent\textbf{题目 \theproblemname 的注记. }}
  {\par}

\begin{document}

\maketitle


\begin{problem}
  证明：(1) $[j^2, l^2]=0$, (2) $[j^2, \boldsymbol{l} \cdot \boldsymbol{s}]=0$，其中 $\boldsymbol{j} = \boldsymbol{l} + \boldsymbol{s}$。
\end{problem}

\begin{solution}
  设总角动量算符
\[
\boldsymbol{j} = \boldsymbol{l} + \boldsymbol{s},
\]
则
\[
j^2 = (\boldsymbol{l}+\boldsymbol{s})^2
= l^2 + s^2 + 2\,\boldsymbol{l}\cdot\boldsymbol{s}.
\]

利用轨道角动量与自旋角动量作用在不同自由度上的事实，有
\[
[l_i,s_j]=0 \quad (\forall\, i,j),
\]
以及角动量的基本对易关系
\[
[l^2,l_i]=0,\qquad [s^2,s_i]=0.
\]

\medskip
\noindent
(1)\;证明 $[j^2,l^2]=0$：

\[
[j^2,l^2]
=[l^2+s^2+2\,\boldsymbol{l}\cdot\boldsymbol{s},\,l^2].
\]
逐项计算：
\[
[l^2,l^2]=0,\qquad [s^2,l^2]=0,
\]
且
\[
[\boldsymbol{l}\cdot\boldsymbol{s},\,l^2]
=\sum_i [l_i s_i,\,l^2]
=\sum_i \bigl(l_i[s_i,l^2]+[l_i,l^2]s_i\bigr)=0.
\]
因此
\[
[j^2,l^2]=0.
\]

\medskip
\noindent
(2)\;证明 $[j^2,\boldsymbol{l}\cdot\boldsymbol{s}]=0$：

\[
[j^2,\boldsymbol{l}\cdot\boldsymbol{s}]
=[l^2,\boldsymbol{l}\cdot\boldsymbol{s}]
+[s^2,\boldsymbol{l}\cdot\boldsymbol{s}]
+2[\boldsymbol{l}\cdot\boldsymbol{s},\boldsymbol{l}\cdot\boldsymbol{s}].
\]
其中
\[
[l^2,\boldsymbol{l}\cdot\boldsymbol{s}]=0,
\qquad
[s^2,\boldsymbol{l}\cdot\boldsymbol{s}]
=\sum_i l_i[s^2,s_i]=0,
\]
并且任意算符与自身的对易子恒为零：
\[
[\boldsymbol{l}\cdot\boldsymbol{s},\boldsymbol{l}\cdot\boldsymbol{s}]=0.
\]
因此
\[
[j^2,\boldsymbol{l}\cdot\boldsymbol{s}]=0.
\]
\end{solution}
\begin{problem}
  求 $\sigma_z$ 在 $|ljm_j\rangle$ 态下的期望值。
\end{problem}

\begin{solution}
要求自旋算符 $\sigma_z$ 在耦合表象 $|l j m_j\rangle$ 下的期望值。

由于
\[
\boldsymbol{\sigma}=2\boldsymbol{s},
\]
只需先计算 $\langle s_z\rangle$。

在角动量耦合态 $|l j m_j\rangle$ 中，有著名结果
\[
\langle l j m_j | s_z | l j m_j \rangle
= \frac{m_j}{2j(j+1)}
\Bigl[ j(j+1)+s(s+1)-l(l+1) \Bigr].
\]
因此
\[
\langle \sigma_z \rangle
=2\langle s_z\rangle
=\frac{m_j}{j(j+1)}
\Bigl[ j(j+1)+s(s+1)-l(l+1) \Bigr].
\]

对于电子自旋 $s=\tfrac12$，有
\[
s(s+1)=\frac{3}{4},
\]
于是
\[
\boxed{
\langle l j m_j | \sigma_z | l j m_j \rangle
=\frac{m_j}{j(j+1)}
\left[
j(j+1)+\frac{3}{4}-l(l+1)
\right]
}.
\]

该结果表明，自旋在 $z$ 方向的期望值与 $m_j$ 成正比，并由轨道角动量 $l$ 与总角动量 $j$ 的耦合关系决定。
\end{solution}


\begin{problem}
  写出 $\boldsymbol{l} \cdot \boldsymbol{s}$ 的本征值和本征态。
\end{problem}

\begin{solution}
考虑轨道角动量 $\boldsymbol{l}$ 与自旋角动量 $\boldsymbol{s}$，
定义总角动量
\[
\boldsymbol{j}=\boldsymbol{l}+\boldsymbol{s}.
\]
由此有恒等式
\[
\boldsymbol{j}^2
=\boldsymbol{l}^2+\boldsymbol{s}^2+2\,\boldsymbol{l}\cdot\boldsymbol{s},
\]
从而
\[
\boldsymbol{l}\cdot\boldsymbol{s}
=\frac{1}{2}\bigl(
\boldsymbol{j}^2-\boldsymbol{l}^2-\boldsymbol{s}^2
\bigr).
\]

由于算符 $\boldsymbol{j}^2,\boldsymbol{l}^2,\boldsymbol{s}^2$ 两两对易，
$\boldsymbol{l}\cdot\boldsymbol{s}$ 的本征态可取为耦合角动量表象
\[
|l s j m_j\rangle.
\]

在该态上，
\[
\boldsymbol{j}^2 |l s j m_j\rangle
= j(j+1)\hbar^2 |l s j m_j\rangle,
\]
\[
\boldsymbol{l}^2 |l s j m_j\rangle
= l(l+1)\hbar^2 |l s j m_j\rangle,
\]
\[
\boldsymbol{s}^2 |l s j m_j\rangle
= s(s+1)\hbar^2 |l s j m_j\rangle.
\]
因此
\[
\boldsymbol{l}\cdot\boldsymbol{s}\,|l s j m_j\rangle
=\frac{\hbar^2}{2}
\Bigl[
j(j+1)-l(l+1)-s(s+1)
\Bigr]
|l s j m_j\rangle.
\]

\medskip
\noindent
综上，$\boldsymbol{l}\cdot\boldsymbol{s}$ 的本征态与本征值为
\[
\boxed{
\begin{aligned}
\text{本征态:}\quad & |l s j m_j\rangle, \\
\text{本征值:}\quad &
\frac{\hbar^2}{2}
\bigl[
j(j+1)-l(l+1)-s(s+1)
\bigr].
\end{aligned}}
\]

对于电子自旋 $s=\tfrac12$，允许的总角动量为
\[
j=l\pm \frac12,
\]
对应的本征值分别为
\[
\boldsymbol{l}\cdot\boldsymbol{s}
=
\begin{cases}
\displaystyle
\frac{\hbar^2}{2}\,l, & j=l+\tfrac12, \\[6pt]
\displaystyle
-\frac{\hbar^2}{2}\,(l+1), & j=l-\tfrac12.
\end{cases}
\]
\end{solution}


\begin{problem}
  电子球面上运动，考虑自旋轨道耦合，其哈密顿量为：
$$H = \frac{\boldsymbol{l}^2}{2\mu R^2} + \gamma \boldsymbol{l} \cdot \boldsymbol{s}$$

(1) 求 $l=0, 1$ 情况下的能级，相应的本征态和简并度；
(2) 若 $t=0$ 时刻电子处于状态 $|l=1, m=1, m_s=-1/2\rangle$，写出 $t>0$ 时刻电子的状态；
(3) 若自旋轨道耦合的方式变为 $\gamma l_x s_x$，求体系所有的能级。
\end{problem}

\begin{solution}
体系哈密顿量为
\[
H=\frac{\boldsymbol{l}^2}{2\mu R^2}+\gamma\,\boldsymbol{l}\cdot\boldsymbol{s},
\qquad s=\frac12 .
\]

利用恒等式
\[
\boldsymbol{l}\cdot\boldsymbol{s}
=\frac12\bigl(\boldsymbol{j}^2-\boldsymbol{l}^2-\boldsymbol{s}^2\bigr),
\]
可得
\[
H=\frac{\boldsymbol{l}^2}{2\mu R^2}
+\frac{\gamma}{2}\bigl(\boldsymbol{j}^2-\boldsymbol{l}^2-\boldsymbol{s}^2\bigr).
\]
因此 $H$ 与 $\boldsymbol{l}^2,\boldsymbol{j}^2,j_z$ 对易，
其本征态可取为 $|l s j m_j\rangle$。

\medskip
\noindent
(1)\;$l=0,1$ 时的能级、本征态与简并度：

\paragraph{$l=0$：}
此时只能有
\[
j=\frac12.
\]
由于 $\boldsymbol{l}\cdot\boldsymbol{s}=0$，
能量为
\[
E_{l=0}
=0.
\]
本征态为
\[
|l=0,s=\tfrac12,j=\tfrac12,m_j=\pm\tfrac12\rangle,
\]
简并度为
\[
g=2.
\]

\paragraph{$l=1$：}
此时
\[
j=\frac32,\;\frac12.
\]

当 $j=\frac32$ 时，
\[
\boldsymbol{l}\cdot\boldsymbol{s}
=\frac12,
\]
能量为
\[
E_{1,3/2}
=\frac{1(1+1)}{2\mu R^2}
+\frac{\gamma}{2}.
\]
简并度
\[
g=2j+1=4.
\]

当 $j=\frac12$ 时，
\[
\boldsymbol{l}\cdot\boldsymbol{s}=-1,
\]
能量为
\[
E_{1,1/2}
=\frac{1(1+1)}{2\mu R^2}
-\gamma.
\]
简并度
\[
g=2j+1=2.
\]

\medskip
\noindent
(2)\;$t=0$ 时处于 $|l=1,m=1,m_s=-\tfrac12\rangle$ 的时间演化：

该态可展开为耦合表象
\[
|1,1,-\tfrac12\rangle
=\sqrt{\frac{1}{3}}\,\Bigl|1,\tfrac32,\tfrac12\Bigr\rangle
-\sqrt{\frac{2}{3}}\,\Bigl|1,\tfrac12,\tfrac12\Bigr\rangle.
\]

对应能量分别为 $E_{1,3/2}$ 与 $E_{1,1/2}$。
因此 $t>0$ 时刻体系状态为
\[
\begin{aligned}
|\psi(t)\rangle
&=\sqrt{\frac13}\,
e^{-iE_{1,3/2}t/\hbar}
\Bigl|1,\tfrac32,\tfrac12\Bigr\rangle
-\sqrt{\frac23}\,
e^{-iE_{1,1/2}t/\hbar}
\Bigl|1,\tfrac12,\tfrac12\Bigr\rangle .
\end{aligned}
\]

\medskip
\noindent
(3)\;自旋轨道耦合改为 $\gamma l_x s_x$ 的情形：

此时哈密顿量为
\[
H=\frac{\boldsymbol{l}^2}{2\mu R^2}+\gamma l_x s_x.
\]

由于
\[
l_x s_x=\frac14
\bigl(l_+s_+ + l_-s_- + l_+s_- + l_-s_+\bigr),
\]
该相互作用在 $|l m\rangle|s m_s\rangle$ 表象中
只耦合有限个态。

对 $l=0$：
\[
l_x s_x=0,
\qquad
E=0,
\]
简并度为 $2$。

对 $l=1$：
在基底
\[
|1,m\rangle|\tfrac12,m_s\rangle
\]
下直接对角化 $l_x s_x$，
可得其本征值为
\[
\lambda=\pm\frac12,\;0.
\]

因此体系总能级为
\[
\boxed{
E=\frac{1(1+1)}{2\mu R^2}
+\gamma\,\lambda,
\qquad
\lambda=\pm\frac12,\;0.
}
\]

这些能级不再由 $j$ 表征，
且简并度由 $l_x s_x$ 的本征值结构决定。
\end{solution}


\begin{problem}
  考虑一个自旋 $s_1 = 1/2$ 和一个自旋 $s_2 = 1$ 组成的体系，
(1) 求总自旋 $S^2, S_z$ 的所有本征值和  本征态（用 $s_{1z}$ 和 $s_{2z}$ 的本征态表示，系数查表）；
(2) 设体系的哈密顿量为 $H = J \boldsymbol{s}_1 \cdot \boldsymbol{s}_2$，求所有的能级及其简并度。
\end{problem}

\begin{solution}
考虑两个自旋：
\[
s_1=\frac12,\qquad s_2=1,
\]
总自旋算符
\[
\boldsymbol{S}=\boldsymbol{s}_1+\boldsymbol{s}_2.
\]

\medskip
\noindent
(1)\;总自旋 $S^2,S_z$ 的本征值与本征态：

由角动量合成规则，
\[
S = s_2+s_1,\; s_2+s_1-1,\;\dots,\;|s_2-s_1|
= \frac32,\;\frac12.
\]

\paragraph{$S=\frac32$ 多重态：}
\[
S^2=\frac32\left(\frac32+1\right)\hbar^2=\frac{15}{4}\hbar^2,
\qquad
M=\frac32,\frac12,-\frac12,-\frac32.
\]

对应的本征态为
\[
\begin{aligned}
\Bigl|\tfrac32,\tfrac32\Bigr\rangle
&=\Bigl|\tfrac12,\tfrac12\Bigr\rangle\,|1,1\rangle,\\[6pt]
\Bigl|\tfrac32,\tfrac12\Bigr\rangle
&=\sqrt{\frac{2}{3}}
\Bigl|\tfrac12,\tfrac12\Bigr\rangle\,|1,0\rangle
+\sqrt{\frac{1}{3}}
\Bigl|\tfrac12,-\tfrac12\Bigr\rangle\,|1,1\rangle,\\[6pt]
\Bigl|\tfrac32,-\tfrac12\Bigr\rangle
&=\sqrt{\frac{2}{3}}
\Bigl|\tfrac12,-\tfrac12\Bigr\rangle\,|1,0\rangle
+\sqrt{\frac{1}{3}}
\Bigl|\tfrac12,\tfrac12\Bigr\rangle\,|1,-1\rangle,\\[6pt]
\Bigl|\tfrac32,-\tfrac32\Bigr\rangle
&=\Bigl|\tfrac12,-\tfrac12\Bigr\rangle\,|1,-1\rangle.
\end{aligned}
\]

\paragraph{$S=\frac12$ 多重态：}
\[
S^2=\frac12\left(\frac12+1\right)\hbar^2=\frac34\hbar^2,
\qquad
M=\frac12,-\frac12.
\]

对应的本征态为
\[
\begin{aligned}
\Bigl|\tfrac12,\tfrac12\Bigr\rangle
&=\sqrt{\frac{1}{3}}
\Bigl|\tfrac12,\tfrac12\Bigr\rangle\,|1,0\rangle
-\sqrt{\frac{2}{3}}
\Bigl|\tfrac12,-\tfrac12\Bigr\rangle\,|1,1\rangle,\\[6pt]
\Bigl|\tfrac12,-\tfrac12\Bigr\rangle
&=\sqrt{\frac{1}{3}}
\Bigl|\tfrac12,-\tfrac12\Bigr\rangle\,|1,0\rangle
-\sqrt{\frac{2}{3}}
\Bigl|\tfrac12,\tfrac12\Bigr\rangle\,|1,-1\rangle.
\end{aligned}
\]

\medskip
\noindent
(2)\;哈密顿量 $H=J\,\boldsymbol{s}_1\cdot\boldsymbol{s}_2$ 的能级与简并度：

利用恒等式
\[
\boldsymbol{s}_1\cdot\boldsymbol{s}_2
=\frac12\bigl(
\boldsymbol{S}^2-\boldsymbol{s}_1^2-\boldsymbol{s}_2^2
\bigr),
\]
其中
\[
\boldsymbol{s}_1^2=\frac34\hbar^2,
\qquad
\boldsymbol{s}_2^2=2\hbar^2.
\]

\paragraph{$S=\frac32$：}
\[
\boldsymbol{s}_1\cdot\boldsymbol{s}_2
=\frac12\left(\frac{15}{4}-\frac34-2\right)\hbar^2
=\frac12\hbar^2,
\]
因此能量为
\[
E_{3/2}=\frac{J}{2}\hbar^2,
\]
简并度
\[
g=2S+1=4.
\]

\paragraph{$S=\frac12$：}
\[
\boldsymbol{s}_1\cdot\boldsymbol{s}_2
=\frac12\left(\frac34-\frac34-2\right)\hbar^2
=-\hbar^2,
\]
因此能量为
\[
E_{1/2}=-J\hbar^2,
\]
简并度
\[
g=2S+1=2.
\]

\medskip
\noindent
综上，体系共有两个能级：
\[
\boxed{
\begin{array}{c|c|c}
S & E & \text{简并度} \\
\hline
\frac32 & \dfrac{J}{2}\hbar^2 & 4 \\
\frac12 & -J\hbar^2 & 2
\end{array}
}
\]
\end{solution}

\end{document}